\chapter{系统分析与设计}

\section{系统概述}

本文设计的企业供应链协同管理系统面向典型制造企业的内部协同与上下游联动场景，围绕订单驱动主线，对销售审核、库存流转、生产执行、采购供应、质量检验与追溯等关键业务进行统一组织。系统采用前后端分离架构，以角色化访问控制为约束，以状态流转为业务驱动，以实时通知和记录留痕为协同支撑，目标是在统一平台内实现跨岗位业务处理和关键数据共享。

从业务定位看，系统并非单一的仓储管理工具或采购登记工具，而是强调订单、库存、生产、采购和质检之间的闭环衔接。顾客发起订单后，系统通过销售审核、仓库处理、生产计划、质检入库和发货执行形成连续流程；原材料短缺时，系统则通过库存预警和采购申请联动采购管理员与供应商，保障生产执行的连续性。由此，系统既关注业务流程效率，也关注过程可追溯与职责边界控制。

\section{需求分析}

\subsection{用户角色分析}

本系统面向多角色协同场景，不同角色的业务职责、可访问资源和操作边界存在显著差异。若系统不能根据角色划分权限范围，容易引发数据越权访问和职责混淆问题。因此，在需求分析阶段需首先明确各类用户角色及其核心职责。

\begin{table}[htbp]
\centering
\footnotesize
\caption{系统主要用户角色及职责}
\label{tab:roles}
\begin{tabular}{|c|p{4.0cm}|p{6.8cm}|}
\hline
角色 & 核心目标 & 主要职责 \\\hline
系统管理员 & 维护系统运行秩序 & 维护用户账号、角色分配、基础数据与权限配置 \\\hline
顾客 & 提交业务需求与查询结果 & 注册登录、下达订单、查看订单处理状态、进行质量追溯查询 \\\hline
销售管理员 & 审核客户订单 & 审核订单有效性、维护销售记录、推动订单进入仓储处理环节 \\\hline
仓库管理员 & 管理库存与出入库 & 维护成品与原材料库存、处理出入库、执行发货、触发库存预警响应 \\\hline
生产管理员 & 组织生产执行 & 接收生产计划、下发生产任务、维护生产记录、协调质检前置环节 \\\hline
采购管理员 & 保障物料供应 & 审核采购申请、生成采购订单、跟踪供应商执行与到货情况 \\\hline
供应商 & 响应采购协同 & 查看采购订单、确认接单、执行发货并反馈发货状态 \\\hline
质检管理员 & 控制质量过程 & 处理来料或成品质检、登记质检结果、支撑质量追溯 \\\hline
\end{tabular}
\end{table}

由表~\ref{tab:roles}可见，本系统的设计难点不在于单一功能点的实现，而在于不同角色围绕同一业务对象执行不同操作时，如何在保证数据共享的同时实现严格的权限边界控制与状态一致性。

\subsection{功能需求分析}

结合业务场景，系统的功能需求主要包括以下几个方面。

第一，用户认证与账号管理需求。系统应支持多角色用户登录、身份校验、角色识别和账号维护，确保不同角色在合法身份前提下访问相应资源。

第二，订单管理需求。顾客能够提交订单，销售管理员完成审核后，订单进入后续仓储、生产和质检处理流程。系统需要完整记录订单创建、审核、执行和交付状态。

第三，库存管理需求。系统需支持成品与原材料库存管理、出入库记录维护、库存数量查询以及库存预警判定，以支撑仓库作业和计划联动。

第四，生产协同需求。当仓库发现成品库存不足以满足订单交付时，系统应能够触发生产计划并将任务传递给生产管理员；生产完成后需进入质检流程，合格后再由仓库入库和发货。

第五，采购供应协同需求。原材料不足时，系统需支持采购申请、采购订单生成、供应商接单发货以及仓库收货入库等操作，从而保证生产连续性。

第六，质检与质量追溯需求。系统应记录关键质量检验信息，并支持顾客或内部人员基于订单信息查询质量处理链路，提升业务透明度。

第七，实时通知与周计划生成功能需求。系统应在关键业务节点主动向相关角色推送消息，并结合库存、安全库存、历史消耗与在途数据生成生产和采购周计划，辅助业务决策。

\subsection{非功能需求分析}

除功能需求外，系统还应满足以下非功能需求。

在安全性方面，系统应具备身份认证、访问控制和敏感操作保护能力，防止越权访问和非法操作。在可用性方面，系统界面应保持操作流程清晰、反馈及时，便于不同岗位快速上手。在一致性方面，系统应确保订单状态、库存数量、生产记录和采购状态在多角色协同处理中保持同步。在可维护性方面，系统应采用清晰的分层结构，便于后续模块扩展和规则调整。在可扩展性方面，系统应支持后续对接移动端、报表分析模块或更细粒度的计划优化模块。

\section{整体架构设计}

\subsection{系统总体架构}

系统总体架构采用典型的前后端分离模式，由前端表示层、后端业务层、数据持久层和消息通信层构成。前端负责提供面向不同角色的业务页面与交互入口；后端负责认证鉴权、流程控制、业务规则校验和数据整合；数据持久层承担结构化业务数据的存储；消息通信层负责关键业务事件的实时通知。

在该架构下，顾客、企业内部各岗位以及供应商均通过浏览器访问系统，前端通过 HTTP 接口与后端进行业务交互，并通过 WebSocket 通道接收实时通知。整体架构既满足多角色统一访问需求，也能够通过服务层将复杂业务规则集中管理，避免将流程判断分散在前端页面中。

\subsection{前后端分层设计}

后端按照控制层、服务层、持久层进行分层设计。控制层负责接收前端请求、进行参数校验并返回标准化响应；服务层封装业务规则、状态流转和事务控制，是系统核心逻辑的承载层；持久层依托 JPA 完成实体管理和数据访问。前端则按照页面层、组件层、接口调用层进行组织，页面层负责业务视图组合，组件层负责通用交互单元复用，接口调用层统一处理请求发送、令牌附带和异常响应。

该分层设计使系统具有较好的职责分离特征。例如，订单审核是否能够进入仓库处理、库存不足时是否触发生产、质检不合格是否允许入库等判断逻辑，均由服务层统一完成，从而保证不同客户端入口下的业务规则一致性。

\subsection{权限控制设计}

权限控制设计遵循“先认证、后授权、按角色分域访问”的原则。系统登录后由后端签发 JWT，后续请求中由安全过滤器解析令牌并恢复用户身份。资源访问控制采用路径级与方法级相结合的方式，对不同角色开放不同接口权限。例如，顾客只允许执行下单、订单查询和追溯查询等操作；销售管理员可审核订单但无权处理库存；仓库管理员可维护库存与发货但无权配置系统账号；供应商仅可访问与自身采购订单相关的业务页面。

此外，权限控制不仅体现在页面是否显示，更体现在服务层业务校验中。即便前端页面进行了入口限制，后端仍需根据当前用户角色和业务状态进行二次校验，从而避免通过构造请求绕过前端控制。

\section{业务流程设计}

\subsection{客户下单与销售审核流程}

顾客在系统中提交订单后，订单首先进入待审核状态。销售管理员对订单内容、交付要求和业务合法性进行审查，审核通过后订单进入仓库处理环节，审核不通过则返回相应处理结果。该流程的设计目的在于保证进入执行阶段的订单具备业务有效性，并由销售角色承担订单入口把关职责。

\subsection{仓储发货与库存流转流程}

仓库管理员接收已审核订单后，首先校验成品库存是否满足发货需求。若库存充足，则执行出库与发货操作，更新库存数量和库存流水；若库存不足，则将缺口信息反馈至生产协同流程。对于已经完成生产并经质检合格的产品，仓库需执行入库操作后再根据订单安排发货。该流程体现了库存作为订单履约关键中间状态的重要作用。

\subsection{生产计划、生产执行与质检流程}

当订单无法直接由现有成品库存满足时，系统根据缺口数量生成生产需求。生产管理员据此制定或接收生产任务，组织生产执行并记录生产过程。生产完成后，产品进入质检环节，由质检管理员记录检验结果。检验合格则允许入库并支持后续发货，不合格则需进入异常处理或返工路径。该流程保证了“先生产、后质检、再入库”的业务顺序，避免未检产品直接进入库存。

\subsection{原材料采购与供应商协同流程}

采购流程既可以由采购管理员主动发起，也可以由库存预警或周计划联动触发。当原材料数量低于安全库存或周计划预测出未来物料缺口时，系统生成采购申请。采购管理员审核后形成采购订单，并向对应供应商发送通知。供应商确认接单后执行发货，仓库在收货后完成入库，必要时可配合质检管理员进行来料质检。该流程强化了企业内部采购部门与外部供应商之间的业务协同。

\subsection{库存预警与周计划生成流程}

库存预警流程通过比较当前库存、在途数量和安全库存阈值，对成品和原材料的潜在缺口进行识别。一旦触发预警，系统向相关角色推送通知，并为周计划生成提供输入依据。周计划生成则综合订单需求、库存余量、历史消耗、在途物料和产品物料清单信息，给出生产周计划与采购周计划建议。该流程的设计目标是将被动响应转化为主动预判。

\subsection{质量追溯流程}

质量追溯流程围绕订单号或产品关联标识展开，将订单信息、生产记录、质检记录和出入库记录进行贯通。当顾客发起质量追溯查询时，系统可展示相应订单在生产、检验和交付过程中的关键处理记录。该流程有助于提升问题定位效率，并增强顾客对企业交付过程透明度的信任。

\section{数据库设计}

\subsection{概念结构设计}

数据库设计围绕业务主对象展开，主要包含用户、角色、产品、订单、订单明细、库存、库存流水、生产计划、生产记录、采购申请、采购订单、供应商、原材料、质检记录和消息通知等实体。各实体通过主外键关系构成完整业务链路。例如，订单与订单明细之间为一对多关系，产品与库存之间存在对应关系，生产计划与订单或库存缺口之间存在关联，采购订单与供应商之间存在从属关系，质检记录则与生产记录或来料记录相关联。

\subsection{主要数据表设计}

\begin{table}[htbp]
\centering
\footnotesize
\caption{系统主要数据表设计}
\label{tab:main-tables}
\begin{tabular}{|c|p{4.0cm}|p{6.8cm}|}
\hline
表名 & 主要字段 & 功能说明 \\\hline
user & id、username、password、status & 存储系统用户基础信息与登录凭据 \\\hline
role & id、role\_name、description & 存储角色定义与权限归属关系 \\\hline
product & id、product\_name、specification、unit & 存储产品基础信息 \\\hline
sales\_order & id、order\_no、customer\_id、status、create\_time & 存储客户订单主信息与状态 \\\hline
sales\_order\_item & id、order\_id、product\_id、quantity & 存储订单明细项 \\\hline
inventory & id、item\_type、item\_id、current\_qty、safe\_qty & 存储成品或原材料库存信息 \\\hline
inventory\_record & id、inventory\_id、change\_type、qty、operator & 存储出入库流水记录 \\\hline
production\_plan & id、plan\_no、product\_id、plan\_qty、status & 存储生产计划信息 \\\hline
purchase\_order & id、purchase\_no、supplier\_id、status、arrival\_status & 存储采购订单与到货状态 \\\hline
quality\_record & id、biz\_type、biz\_id、result、remark & 存储质检结果与追溯依据 \\\hline
message\_notice & id、receiver\_id、message\_type、is\_read & 存储实时通知留痕记录 \\\hline
\end{tabular}
\end{table}

表~\ref{tab:main-tables}给出了系统主要数据表设计。通过上述表结构，系统能够对订单执行、库存变动、生产任务、采购协同、质检记录和消息通知进行统一管理。

\subsection{关键实体关系分析}

在关键实体关系中，用户与角色之间为多对多关系，用于支持同一用户持有一个或多个角色；订单与订单明细之间为一对多关系，确保一个订单可以包含多个产品项；产品与库存之间通过业务类型区分成品库存与原材料库存；生产计划通常关联特定产品和触发来源；采购订单与供应商之间为多对一关系；质检记录通过业务标识与生产或来料过程相关联；消息通知则与接收用户和业务事件绑定。

实体关系设计的核心目标在于支撑系统业务闭环。一方面，各实体需能独立维护自身状态和属性；另一方面，不同实体之间必须通过关联关系反映真实业务链路，以保证查询、追溯和状态传播的正确性。这一设计为后续详细实现提供了数据层基础。